\documentclass[12pt,a4paper]{article}

% ----- Encoding & language -----
\usepackage[utf8]{inputenc}
\usepackage[T1]{fontenc}
% NOTE: babel[spanish] requires texlive-lang-spanish, absent on some build
% machines; omitted so the report compiles anywhere. Content is Spanish; only
% hyphenation falls back to default. Re-add \usepackage[spanish]{babel} where
% the package is installed for proper Spanish hyphenation.

% ----- Math -----
\usepackage{amsmath,amssymb}

% ----- Layout & typography -----
\usepackage[margin=2.5cm]{geometry}
\usepackage{microtype}
\usepackage{parskip}
\usepackage{setspace}
\onehalfspacing

% ----- Tables -----
\usepackage{booktabs}
\usepackage{multirow}
\usepackage{array}
\usepackage{xcolor}

% ----- Graphics & floats -----
\usepackage{graphicx}
\usepackage{float}
\usepackage{caption}

% ----- Hyperlinks -----
\usepackage[colorlinks=true, linkcolor=blue!60!black, citecolor=blue!60!black, urlcolor=blue]{hyperref}

% ----- Misc -----
\usepackage{enumitem}

% ================================================================
\title{\textbf{Informe de Avances}\\[0.4em]
       \large Predicción de Irradiancia Solar (GHI) con Aprendizaje Profundo\\
       y Cuantificación de Incertidumbre}
\author{Esteban Ladino Nieto}
\date{5 de agosto de 2026}
% ================================================================

\begin{document}
\maketitle

% ================================================================
\section{Contexto}
% ================================================================

El objetivo del trabajo es pronosticar la irradiancia solar horizontal
global (GHI) a 1, 3 y 6~horas usando imágenes del satélite GOES-16 como
entrada espacial, en dos sitios colombianos con regímenes climáticos
contrastantes (El Paso, César --- tropical seco de baja elevación; Uniandes,
Bogotá --- montaña ecuatorial con convección orográfica). El sistema además
debe producir una capa de cuantificación de incertidumbre \emph{epistémica}
(inicialmente vía SGLD).

Este informe cubre varias semanas densas en incidentes. Se organiza por
problema encontrado, en orden de gravedad para la validez de los resultados:
(i)~un bug de datos que invalidó silenciosamente todo El Paso;
(ii)~un error metodológico en el baseline SARIMA; (iii)~problemas de
infraestructura de cómputo (errores de CUDA, reinicio y actualización del
servidor); y (iv)~un problema de fondo, aún abierto, en el método de
cuantificación de incertidumbre. Al final se resume el estado actual y lo
pendiente. La política de trabajo adoptada en todo momento fue: \emph{ante
datos o métodos con error, marcar y detener, nunca publicar un número que no
se sostenga}.

% ================================================================
\section{Incidente 1: desalineación temporal en El Paso}
\label{sec:incidente}
% ================================================================

\subsection{Qué se encontró}

El archivo de referencia terrestre de El Paso
(\texttt{ground\_10min\_utc\_elpaso.parquet}) tenía \textbf{todos sus
timestamps de GHI desfasados 5~horas} respecto a la hora UTC real. El pico
diario de GHI, que debería caer alrededor de las 16--17~UTC (mediodía solar
para esa longitud), aparecía en la hora 11~UTC.

\subsection{Causa raíz}

El CSV crudo de origen (\texttt{data\_raw/Data\_ElPaso\_tierra.csv}) trae un
sufijo \texttt{+00:00} escrito directamente en cada timestamp, aunque los
valores son en realidad hora local de Bogotá. El notebook
\texttt{01\_ground\_eda.ipynb} lo heredó sin corregirlo: la rama de Uniandes
sí localiza correctamente a \texttt{America/Bogota} antes de convertir a
UTC; la rama de El Paso, al encontrar el timestamp ya marcado como UTC
(aunque incorrectamente), se saltaba ese paso.

Como \texttt{04\_build\_manifests.py} empareja cada secuencia satelital con
su GHI objetivo por coincidencia exacta de timestamp contra el índice de
cobertura satelital (que sí está en UTC real), el desfase de 5~horas hacía
que cada muestra de El Paso emparejara imágenes satelitales de un momento
real distinto al de su etiqueta de GHI --- de forma silenciosa, sin generar
ningún error ni muestra faltante.

\subsection{Impacto y corrección}

\textbf{Solo El Paso se vio afectado.} Uniandes estaba correcto desde el
principio (su pico de GHI cae en la hora~16~UTC). Persistencia de El Paso
depende de la consistencia interna de la propia serie terrestre, así que
también cambió al corregir el archivo (204.4~$\to$~206.0~W/m² a 1h); las
estadísticas puramente descriptivas de forma/distribución recomputadas hoy
sobre el archivo corregido son válidas. Pero \textbf{todo modelo de
aprendizaje profundo entrenado con El Paso} --- baseline, Optuna~v1 (los
resultados insignia que se citaban, incluido $\text{Skill}_\text{day}=0.636$
en 6h), Optuna~v2, MLP, fusión, SGLD --- estaba construido sobre el manifest
desalineado y debía tratarse como inválido.

La corrección: se arregló el notebook y se regeneró el archivo de referencia
(pico ahora en 16~UTC, verificado); se reconstruyó el manifest de El Paso
para los 3~horizontes y se completó el patch store; se recalcularon
persistencia y SARIMA sobre datos corregidos; se \textbf{borraron todos los
runs de El Paso} construidos sobre el manifest viejo (Uniandes intacto); y se
relanzó el entrenamiento de El Paso desde cero. Un detalle operativo
importante: \texttt{data/} está en \texttt{.gitignore}, así que la corrección
no viajó al servidor por git y hubo que sincronizarla por \texttt{rsync}
explícitamente; el primer relanzamiento en el servidor corrió $\sim$40~min
sobre el manifest aún viejo antes de detectarlo.

\subsection{Cómo se refleja en el artículo (convención de la daga)}

Semanas después se descubrió que el manuscrito \emph{todavía} citaba los
números contaminados de El Paso en tablas, figuras y prosa. En vez de borrar
o reestructurar esas secciones, se adoptó una convención de marcado: cada
valor derivado de datos pre-corrección lleva una \textbf{daga}
($^\dagger$) y una nota de ``pendiente de reconfirmar''; secciones,
hipótesis y figuras se conservan íntegras. La figura de serie temporal de El
Paso se dejó como recuadro-marcador (``pendiente de regeneración''); las
figuras de barras muestran El Paso con solo las barras actualmente válidas
(persistencia y SARIMA), sin ocultar el sitio.

Esta decisión resultó crucial: el primer resultado corregido de El Paso
(ResNet a 1h, ya con sus 2~semillas) da $\text{Skill}_\text{day} = +0.297
\pm 0.028$, cuando el número contaminado era $-0.167^\dagger$
(\emph{negativo}). El bug no desviaba el número: le \emph{invertía el signo}.
La hipótesis actual del artículo de que ``el horizonte de 1h en El Paso es
intratable'' tendrá que reescribirse, no solo actualizarse, cuando se
des-dague esa celda. Haberlo marcado como pendiente evitó publicar una
conclusión falsa.

% ================================================================
\section{Incidente 2: error metodológico en SARIMA}
% ================================================================

El baseline SARIMA producía un pronóstico estático: ajustaba el modelo una
vez y luego generaba \emph{una sola trayectoria abierta} desde el fin del
entrenamiento que cubría meses de validación y test. En consecuencia sus
``pronósticos a 1h, 3h y 6h'' eran idénticos entre horizontes (solo cambiaba
el subconjunto de test evaluado), y cada objetivo de test era en realidad una
extrapolación a miles de pasos. Como benchmark por horizonte, era inválido.

Se reescribió a \textbf{evaluación rolling-origin}: parámetros estimados una
vez sobre el entrenamiento y luego fijos, con el filtro de Kalman avanzando
observación por observación por validación y test, emitiendo en cada origen
horario de test un pronóstico genuino a $H$~pasos desde el estado que
incorpora la información hasta ese instante. Los nuevos valores hacen a
SARIMA un baseline \emph{competitivo}: $\text{Skill}_\text{day}$ de
$0.42/0.65/0.74$ (El Paso) y $0.33/0.41/0.48$ (Uniandes) a 1h/3h/6h. Se
documentó explícitamente en la metodología que estos números \emph{no} son
comparables directos con la tabla principal, porque SARIMA se evalúa sobre el
subconjunto horario y contra un objetivo de media horaria (más suave que el
instantáneo de 10~min que resuelven los modelos de aprendizaje profundo).

% ================================================================
\section{Incidente 3: infraestructura de cómputo (CUDA, reinicio, versiones)}
% ================================================================

\subsection{Error de CUDA por \texttt{fork} de workers}

Al agregar procesos concurrentes en el servidor empezaron a aparecer fallos
\texttt{CUDA error: initialization error}: el problema conocido de PyTorch en
que un worker de \texttt{DataLoader} creado por \texttt{fork()} hereda un
contexto CUDA ya inicializado por el proceso padre. Inicialmente pareció que
2~procesos concurrentes eran estables, pero el fallo terminó apareciendo
incluso con 2~vías tras acumular suficientes \texttt{fork}s de trials. La
causa no es la concurrencia sino que la condición de carrera está presente en
\emph{cada} trial (cada uno crea workers nuevos en un proceso cuyo contexto
CUDA lleva vivo desde el trial~1). \textbf{Primer arreglo}: se puso
\texttt{num\_workers=0} en todos los scripts de entrenamiento y en SGLD.
Sin \texttt{fork} no hay contexto heredado, así que la carrera desaparece de
raíz.

\subsection{Optimización: carga de datos paralela con \texttt{spawn}}

El arreglo \texttt{num\_workers=0} era seguro pero dejaba la carga de datos en
un solo hilo, así que los trabajos quedaban limitados por CPU y \textbf{la GPU
trabajaba al 13--21\%} con la mayoría de los núcleos ociosos. Se resolvió de
raíz cambiando el \emph{start method} de los workers de \texttt{fork} a
\texttt{spawn}: los procesos \texttt{spawn} arrancan sin heredar contexto CUDA
(la carrera desaparece igual) y permiten \texttt{num\_workers>0} sin riesgo.
Medido: \textbf{$\sim$2$\times$ más rápido por batch} y, en la práctica, la
\textbf{GPU pasó de $\sim$20\% a 99\% de uso} (los trabajos ya no la dejan
esperando), sin costo de memoria GPU (los workers solo alimentan el modelo).
Se validó de punta a punta antes de aplicarlo y se propagó a ambas máquinas.
Un detalle atrapado por la prueba end-to-end (no por el microbenchmark): bajo
\texttt{spawn} la función de \emph{seeding} de workers debía ser picklable, lo
que exigió moverla a nivel de módulo.

\subsection{Desincronización de driver en la máquina local}

En la máquina local, \texttt{nvidia-smi} empezó a fallar con
\emph{driver/library version mismatch} (el driver se actualizó en disco sin
reiniciar). Se verificó que esto afecta solo a la interfaz de monitoreo
(NVML): un proceso CUDA nuevo inicializa y corre sin problema, y los jobs ya
vivos no se ven afectados. La única precaución operativa fue no reiniciar la
máquina hasta que terminaran los trabajos en curso.

\subsection{Reinicio y actualización del servidor}

El servidor del asesor estuvo inalcanzable (pérdida total de red) y al volver
había \textbf{reiniciado} --- por \texttt{unattended-upgrades} del sistema.
Los runs en curso murieron (los estudios Optuna pierden su historial en
memoria); lo ya completado sobrevivió en disco. El asesor confirmó que
aprovechó el reinicio para actualizar el entorno. Se hizo un \emph{sanity
check} antes de relanzar: el stack pasó de \texttt{torch~2.11} a
\texttt{torch~2.14} (\texttt{pandas~3.0}, \texttt{numpy~2.4} y
\texttt{optuna~4.7} coinciden con el local). Se comprobó con un smoke-test
completo que el pipeline funciona en el stack nuevo y que los checkpoints
viejos (entrenados en 2.11) cargan bien en 2.14. Queda como caveat menor de
reproducibilidad que El Paso quedará con una mezcla de \texttt{torch~2.11}
(lo completado antes del reinicio y lo que corre en local) y \texttt{2.14}
(lo relanzado en el servidor), diferencia dentro del ruido entre semillas ya
que se reporta media$\pm$desviación.

\textbf{Costo real: reinicios inesperados en máquinas compartidas.} Tanto el
servidor como la máquina local son equipos compartidos, y ambos han
reiniciado sin aviso (por \texttt{unattended-upgrades} u otros usuarios). Cada
reinicio mata los estudios Optuna en curso, que pierden su historial en
memoria y deben relanzarse desde cero. Un reinicio local costó
$\sim$6.5~días de cómputo de El Paso. Esto es la principal fuente de
resbalón en el cronograma --- más que la velocidad de cómputo --- y conviene
coordinar los reinicios para no perder semanas.

% ================================================================
\section{Cuantificación de incertidumbre: objetivo, avance y problema abierto}
\label{sec:sgld}
% ================================================================

\subsection{Qué quiero hacer: desentrelazar incertidumbre aleatoria y epistémica}

La meta no es solo dar un intervalo alrededor del pronóstico, sino
\textbf{descomponer} la incertidumbre en sus dos fuentes, porque la respuesta
operativa de un operador de red es distinta en cada caso:

\begin{itemize}[leftmargin=1.5em, noitemsep]
  \item \textbf{Incertidumbre aleatoria} $\sigma_a^2(x)$: el ruido
        \emph{irreducible} del proceso físico --- la estocasticidad de la
        formación de nubes. Aunque tuviéramos datos infinitos, seguiría ahí.
        Ante ella, la respuesta es una reserva o un margen conservador.
  \item \textbf{Incertidumbre epistémica} $\sigma_\text{epi}^2(x)$: la
        incertidumbre \emph{reducible} del modelo, alta cuando pronostica
        bajo condiciones que ha visto poco (datos escasos o atípicos). Ante
        ella, la respuesta es distinta: recolectar más datos de esa condición
        o reentrenar, no solo ampliar el margen.
\end{itemize}

El método que sigo es el propuesto por el asesor: la estrategia
\emph{cooperativa} de Yi \& Bessa (2025), \emph{Cooperative Bayesian and
Variance Networks Disentangle Aleatoric and Epistemic Uncertainties}
(arXiv:2505.02743). Su idea central es entrenar las componentes en
\textbf{secuencia, no en conjunto}: entrenar la media y la varianza a la vez
(las redes MVE estándar) produce gradientes desbalanceados y varianza
sobreconfiada. El método tiene tres pasos:

\begin{enumerate}[leftmargin=1.6em, noitemsep]
  \item \textbf{Red de media} --- el backbone de punto ya entrenado (Optuna),
        congelado.
  \item \textbf{Red de varianza (aleatoria)} --- con la media congelada,
        entrena una red que predice el residuo al cuadrado
        $r=(\mu(x)-y)^2$ vía verosimilitud Gamma, dando
        $\sigma_a^2(x)=\alpha(x)/\lambda(x)$.
  \item \textbf{Red bayesiana (epistémica)} --- muestrea la posterior de los
        pesos (SGLD), \emph{arrancando desde la media} y usando una
        verosimilitud ponderada por la $\sigma_a^2(x)$ fija del paso~2, para
        que la dispersión resultante sea epistémica genuina, ya separada del
        ruido aleatorio.
\end{enumerate}

\subsection{Lo que he hecho: pasos 1 y 2}

\textbf{Paso 1 (media):} listo --- son los backbones de Optuna por
sitio-horizonte, con sus checkpoints guardados.

\textbf{Paso 2 (varianza aleatoria):} \textbf{implementado y validado}. Se
programó la red de varianza con un truco limpio (reemplazar la última capa
\texttt{Linear($\cdot$,1)} del backbone por \texttt{Linear($\cdot$,2)} para
que emita los dos parámetros Gamma sin tocar el \texttt{forward} original) y
la verosimilitud Gamma \emph{derivada desde cero} --- no transcrita de la
Ec.~6 del paper, que se leyó de una imagen del PDF con posible error de OCR.
La validación fue triple: (i)~recuperación sintética de parámetros Gamma
conocidos; (ii)~validación cruzada (el residuo medio del test coincide con el
\texttt{mae\_day} del checkpoint original, 199.51 vs 199.31~W/m²); y
(iii)~estabilidad (15--30 épocas sin divergencia en las 3 arquitecturas). Se
corrió sobre Uniandes (único sitio con backbones disponibles tras la
reconstrucción de El Paso): 33/36 combinaciones (ResNet 12/12, GraphSAGE
12/12, MLP 9/12).

\emph{Evidencia de que la parte aleatoria funciona.} Midiendo la correlación
entre la $\sigma_a$ estimada y el error real $|\mu(x)-y|$ (una $\sigma_a$
útil debe ser \emph{heteroscedástica}, es decir, mayor donde el error es
mayor), la mayoría de combinaciones dan estimaciones significativas: ResNet a
1h alcanza correlación $0.53$, y GraphSAGE/ResNet a 3h y 6h están en
$0.33$--$0.41$. \emph{El colapso homoscedástico se concentra a 1h} (GraphSAGE
y FlatMLP colapsan a varianza casi constante en 3 de 4 semillas), un modo de
fallo conocido de las redes de varianza, coherente porque a ese horizonte el
residuo del modelo de punto tiene poca señal explotable (su skill también es
el más débil ahí). Conclusión operativa: la mitad aleatoria del
desentrelazamiento \textbf{ya funciona} para reportar, exceptuando o notando
el caso de 1h; conviene revisar ese colapso (¿lr o épocas distintas para la
red de varianza?) antes de los resultados finales.

\subsection{El problema abierto: paso 3 (red bayesiana / SGLD)}

El paso~3 es donde el proyecto se atascó. La red bayesiana es el SGLD ya
existente, pero requería adaptarse (warm-start desde la media +
verosimilitud ponderada por $\sigma_a^2$). Al intentar validar SGLD sobre
datos de resolución completa --- algo que \emph{nunca se había hecho}, sus
hiperparámetros venían de un montaje anterior mucho más barato --- falló en
cascada. En orden (todo verificado empíricamente sobre Uniandes):

\begin{enumerate}[leftmargin=1.6em, noitemsep]
  \item \textbf{Presupuesto irreal.} El programa original (500~épocas de
        burn-in + 10~muestras cada 100 = 1500~épocas) asumía un costo por
        época $\sim$30$\times$ menor del real ($\sim$450~s/época sobre 54k
        muestras): implicaba $\sim$8~días por corrida.
  \item \textbf{Init fresco no entrena.} Reducido el programa, con
        inicialización aleatoria la red no aprende: al paso de muestreo
        ($\varepsilon = 10^{-5}$) los updates tipo SGD no alcanzan a entrenar
        desde cero; la pérdida se queda en el piso de ``predecir la media''.
  \item \textbf{Warm-start + prior fuerte destruye el ajuste.} Arrancar desde
        el checkpoint entrenado carga bien, pero el prior gaussiano centrado
        en cero con precisión $\lambda=100$ (elegido para confinar la cadena)
        arrastra los pesos entrenados de vuelta al origen y degrada el
        ajuste.
  \item \textbf{Warm-start + prior débil + paso grande diverge.} Con un prior
        débil, a $\varepsilon = 10^{-5}$ la cadena diverge en pocas épocas
        (pérdida $0.9\to78$): cada época son $\sim$3400~pasos SGLD, demasiado
        ruido acumulado.
  \item \textbf{El paso calibrado parecía funcionar\ldots\ y era un falso
        positivo.} Un barrido dio $\varepsilon = 10^{-7}$ como estable en una
        ventana corta (4~épocas): pérdida estable, skill del ensamble
        $0.184$ ($\geq$ backbone $0.175$), $\sigma_\text{epi}$ razonable.
        Pero sobre las 70~épocas reales la cadena \textbf{deriva lentamente}:
        estable $\sim$10~épocas y luego la pérdida crece ($0.4\to1.7$) y el
        val se degrada; los $\sigma_\text{epi}$ finales explotan
        ($\sim$120--200 vs 17 en la ventana estable). Las 3~corridas dieron
        skill negativo o degradado.
\end{enumerate}

\textbf{Diagnóstico.} El problema es fundamental, no de ajuste fino: SGLD
sobre redes sobreparametrizadas no tiene una posterior estacionaria acotada
--- la cadena hace un paseo aleatorio por las direcciones planas de la
función de pérdida y se aleja del modo. El prior lo bastante fuerte para
acotar esa deriva es también lo bastante fuerte para distorsionar la
solución. Es una tensión conocida y difícil en la literatura (``cold
posterior''). Todo esto quedó documentado en la metodología del artículo como
\emph{problema abierto}, sin afirmar que el estimador esté validado, y los
runs fallidos se archivaron como evidencia.

\subsection{Lo que sigue}

\begin{enumerate}[leftmargin=1.6em, noitemsep]
  \item \textbf{Completar el paso~2} en las combinaciones de Uniandes que
        faltan, y en El Paso una vez existan sus backbones corregidos.
  \item \textbf{Definir la componente epistémica (paso~3) con el asesor.} Dado
        que el SGLD tal cual no da una posterior estable, hay dos caminos y la
        decisión es suya, ya que el método bayesiano fue su propuesta:
        \begin{itemize}[leftmargin=1.3em, noitemsep]
          \item \emph{Arreglar/reemplazar el muestreador} bayesiano (p.\,ej.\
                un prior intermedio que acote las direcciones planas sin
                distorsionar el modo, o un esquema tipo SGHMC con momento) ---
                trabajo de investigación de resultado incierto.
          \item \emph{Sustituir la red bayesiana por un deep ensemble}: usar
                la dispersión de las 4~semillas ya entrenadas por
                configuración como estimador epistémico
                (Lakshminarayanan et~al., 2017), un estándar bien calibrado.
                No requiere entrenamiento nuevo y \emph{preserva el objetivo
                de desentrelazamiento}: el paso~2 (red de varianza) seguiría
                dando la parte aleatoria $\sigma_a^2$, y el ensemble daría la
                epistémica $\sigma_\text{epi}^2$ --- la misma descomposición
                que buscaba Yi \& Bessa, con una componente epistémica más
                robusta.
        \end{itemize}
  \item \textbf{Revisar el colapso homoscedástico} del paso~2 en GraphSAGE~a~1h
        antes de reportar resultados finales.
\end{enumerate}

Mientras se decide, la Tabla de descomposición de incertidumbre queda como
marcador declarado en el artículo, sin comprometer un método que no esté
validado.

% ================================================================
\section{Estado actual}
\label{sec:estado}
% ================================================================

\begin{table}[H]
  \centering
  \small
  \caption{Estado por grupo de experimento al 5 de agosto de 2026.
           ``v2'' = espacio extendido, 2~semillas $\{42,1\}$, 75~trials.
           Uniandes intacto desde el incidente; El Paso reconstruyéndose en
           paralelo (servidor + local).}
  \label{tab:estado}
  \footnotesize
  \begin{tabular}{@{}lp{4.2cm}cp{3.9cm}@{}}
    \toprule
    \textbf{Grupo} & \textbf{El Paso} & \textbf{Uniandes} & \textbf{Observación} \\
    \midrule
    Optuna v1 (ResNet/GSAGE) & ---  & 24/24 & Base Tabla~2 Uniandes (válida) \\
    SARIMA + persistencia    & \multicolumn{2}{c}{ok} & Reescrito a rolling-origin \\
    \midrule
    ResNet Optuna v2    & h1 ok; h3, h6 corren & casi & Servidor + local \\
    GraphSAGE Optuna v2 & h1, h3 (s42 ok); resto corre & casi & Servidor + local \\
    FlatMLP Optuna      & 1 semilla & completo & \textbf{Uniandes ya en Tabla~2} \\
    Fusion ResNet+LSTM  & h6 ok & 2/6 & En cola el resto \\
    Fusion GraphSAGE    & pendiente & pendiente & En cola \\
    Red de varianza     & --- & 33/36 & Paso 2 UQ validado \\
    SGLD (epistémico)   & \multicolumn{2}{c}{en revisión} & Ver Sección~\ref{sec:sgld} \\
    \bottomrule
  \end{tabular}
\end{table}

\textbf{Resultados corregidos de El Paso (fortalecen el paper).} Los primeros
números con datos corregidos superan a los contaminados: GraphSAGE~v2 a 3h da
$\text{Skill}_\text{day} = 0.585$ (vs $0.419^\dagger$), y ResNet~a~1h da
$+0.297 \pm 0.028$ (vs $-0.167^\dagger$, que era negativo). La corrección no
solo no daña El Paso: lo mejora. \textbf{La columna de 1h está por completarse}
--- ResNet ya tiene sus 2 semillas y GraphSAGE~a~1h va en $\sim$43/75 trials
de su segunda semilla; al terminar será la primera columna des-dagueable, y
obligará a reescribir la hipótesis de que ``1h es intratable''.

\textbf{Reparto de cómputo y eficiencia.} Servidor (cadena
\texttt{run\_sequential}): El Paso~v2 sistemático (h1, h3). Local: El Paso~a~6h
(la cola de la cadena, para paralelizar sin duplicar). Tras la optimización de
\texttt{spawn}, ambas máquinas ahora saturan su GPU ($\sim$99\% local) en vez
de dejarla ociosa; el cuello de botella pasó a ser la memoria GPU (los trials
de GraphSAGE con vecindario y ancho grandes usan hasta $\sim$9~GB), no la CPU,
lo que limita cuántos trabajos caben en paralelo.

% ================================================================
\section{Progreso del artículo}
% ================================================================

Se hizo una revisión crítica del manuscrito desde la perspectiva de un
revisor de revista y se aplicaron las correcciones que no dependen de
experimentos pendientes:

\begin{itemize}[leftmargin=1.5em, noitemsep]
  \item Convención de la daga aplicada en todo El Paso (Sección~\ref{sec:incidente}):
        el artículo es auto-consistente y honesto sobre qué está pendiente.
  \item Metodología puesta al día con la implementación real: protocolo
        Optuna~v2 (75~trials, 2~semillas), SARIMA rolling-origin, y el
        muestreador de incertidumbre (marcado en revisión, no afirmado como
        validado).
  \item Correcciones editoriales previas: error numérico en el resaltado de
        la Tabla~2, sobre-afirmación sobre SGLD atenuada, gap de motivación
        para UQ agregado a la introducción, transiciones en Trabajo
        Relacionado, y 4~figuras descriptivas del dataset integradas.
  \item Extensión cooperativa de UQ (red de varianza aleatoria) descrita y
        con su paso~2 implementado.
  \item \textbf{Fila de FlatMLP (ablación espacial) llenada para Uniandes} en
        la Tabla~2, con un resultado que \emph{refuerza la narrativa del
        paper}: FlatMLP (que hace \emph{average-pool} del parche y descarta
        toda estructura espacial) es competitivo con los modelos espaciales a
        horizontes largos en Uniandes ---a 6h da $\text{Skill}_\text{day} =
        0.401$, por encima de ResNet ($0.368$) y casi igual a GraphSAGE
        ($0.417$)--- y solo se abre una brecha clara a 1h. Es evidencia
        directa de que en este sitio de montaña tropical las features
        espaciales del satélite aportan poco valor marginal, justo lo que la
        discusión hipotetiza.
\end{itemize}

Uniandes queda con resultados finales y defendibles (4~semillas, auditado
de punta a punta: sin fugas de información, normalizador consistente,
persistencia evaluada sobre las mismas muestras que los modelos). El Paso
queda honestamente marcado como pendiente. Las tablas que dependen de
experimentos en curso (El Paso, ablación $k$-NN, descomposición de
incertidumbre) se llenarán conforme terminen.

% ================================================================
\section{Decisiones}
% ================================================================

\textbf{Tomadas} (ya aplicadas en el código y/o el artículo):

\begin{itemize}[leftmargin=1.5em, noitemsep]
  \item \textbf{Convención de la daga} para datos con error: marcar y notar
        como pendiente, nunca borrar secciones/hipótesis/figuras. Evitó
        publicar una conclusión falsa (el caso de 1h en El Paso).
  \item \textbf{SARIMA a rolling-origin}, y documentar que no es comparable
        directo con la tabla principal (subconjunto horario, objetivo de
        media horaria).
  \item \textbf{El Paso usará solo el grafo $k$-NN tuneado} (v2) para
        GraphSAGE; \emph{no} se resucita el código del grafo fijo ya
        discontinuado --- El Paso no tendrá su propia ablación fijo-vs-tuneado.
  \item \textbf{Carga de datos con \texttt{spawn}} en vez de
        \texttt{num\_workers=0} (2$\times$, GPU al 99\%).
  \item \textbf{SGLD marcado como problema abierto} en la metodología del
        artículo, sin afirmarlo como validado; Tabla de descomposición como
        marcador declarado.
\end{itemize}

\textbf{Por tomar} (idealmente en la reunión con el asesor):

\begin{enumerate}[leftmargin=1.6em, noitemsep]
  \item \textbf{Método de incertidumbre epistémica (la decisión principal).}
        Como el SGLD no da una posterior estable, elegir entre: (a)~invertir
        en arreglar el muestreador bayesiano (resultado incierto), o
        (b)~usar \emph{deep ensembles} sobre las 4~semillas ya entrenadas
        (inmediato, robusto, preserva el desentrelazamiento con el paso~2).
        Esta decisión desbloquea la Tabla de descomposición.
  \item \textbf{Colapso homoscedástico de la red de varianza a 1h}: ¿reajustar
        (lr/épocas) o reportarlo como limitación?
  \item \textbf{Comparación SARIMA vs.\ aprendizaje profundo en base horaria}
        (posible pedido de un revisor): ¿hacerla antes de someter o dejarla
        como trabajo futuro declarado?
  \item \textbf{Infraestructura}: coordinar reinicios de las máquinas
        compartidas para no perder cómputo; reconciliar el historial de git
        divergente del servidor (pendiente menor).
\end{enumerate}

% ================================================================
\section{Trabajo pendiente --- priorizado}
% ================================================================

\begin{enumerate}[leftmargin=1.5em, noitemsep]
  \item Dejar correr El Paso~v2 en servidor (h1, h3) y local (h6). Las
        corridas ya completadas en el servidor se sincronizaron al local
        (por \texttt{rsync}, solo los \texttt{summary.json}) para que la
        agregación las vea.
  \item \textbf{Des-daguear cada columna de El Paso} conforme tenga sus
        2~arquitecturas $\times$ 2~semillas, reescribiendo la narrativa
        asociada. La columna de 1h es la más cercana (falta una semilla) y
        será la primera; su signo cambió, así que implica reescribir, no solo
        actualizar.
  \item Completar \texttt{fusion} y el FlatMLP de El Paso (en cola).
  \item Aplicar la decisión de incertidumbre epistémica (Sección~Decisiones)
        y llenar la Tabla de descomposición.
  \item Ejecutar las decisiones de infraestructura (coordinar reinicios;
        reconciliar el git del servidor).
\end{enumerate}

% ================================================================
\section{Conclusión}
% ================================================================

El periodo estuvo dominado por la depuración de errores que ponían en riesgo
la validez de los resultados: un bug de datos que invertía el signo del
desempeño de El Paso, un baseline SARIMA mal planteado, fallos de CUDA por
\texttt{fork}, un reinicio con actualización del servidor, y --- lo más
significativo y aún abierto --- que el método de incertidumbre epistémica
elegido (SGLD) no produce una posterior estable sobre estas redes. Cada uno
se detectó, aisló y documentó antes de que un número inválido llegara al
artículo. La parte del trabajo que sí es sólida (Uniandes, SARIMA
rolling-origin, el pipeline auditado, El Paso honestamente marcado) está
intacta; el riesgo restante es de cronograma (El Paso toma semanas) y de una
decisión metodológica pendiente sobre la cuantificación de incertidumbre.

\end{document}
